\clearpage
\fancyhead[L]{\small MONDSTERNE.}
\hypertarget{native-448-mondsterne}{}
\begin{center}
{\large MONDSTERNE.}
\end{center}

\begin{center}
Astronomische Nachrichten. Band II. Nr. 26. Seite 17..20. 1823 Februar.
\end{center}

\begin{center}
Berechnung der Meridiandifferenz zweier Orte aus entsprechenden Mondkulminationen.\\
Von Herrn Professor \textsc{Nicolai}, Direktor der Sternwarte in Mannheim [an \textsc{Schumacher.}].
\end{center}

Mit diesen Zeilen will ich heute versuchen, Ihrer freundlichen Bitte nachzukommen und Ihnen mitzuteilen, wie ich aus den entsprechenden Mondkulminationen die Meridiandifferenz berechne.

Diese Methode beruht bekanntlich darauf, dass man an zwei Orten den wahren Unterschied der Rektaszension zwischen einem der Mondränder und einigen zuvor vereinbarten Fixsternen beobachtet. Der Unterschied zwischen der wahren Rektaszension des Mondrandes und der eines Fixsterns sei (positiv gerechnet, wenn die Rektaszension des Mondes größer ist als die des Sterns) am westlichen Beobachtungsort $=t$ und am östlichen $=\tau$. Dann ist $t-\tau$ der wahre Rektaszensionsunterschied des Mondrandes für die zwischen beiden Beobachtungen verstrichene Zeit. Hat man an beiden Orten mehrere Fixsterne verglichen, so nimmt man das arithmetische Mittel der verschiedenen Werte von $t-\tau$, sofern die Beobachtungen an beiden Orten gleich genau waren. Wie vorzugehen ist, wenn das hinsichtlich der Zahl der beobachteten Fäden nicht der Fall ist, werden wir weiter unten sehen. --- Offenbar kommt es aber darauf an, den wahren Rektaszensionsunterschied des \emph{Mittelpunkts} des Mondes zu erhalten. Dieser wird von dem eben erhaltenen etwas abweichen, weil sich der Halbmesser des Mondes in AR. zwischen den Beobachtungen ändert, teils durch die Änderung des wahren Mondhalbmessers selbst, teils durch die Änderung der Deklination des Mondes. Bezeichnet man für die Kulminationszeit des Mondes am westlichen Beobachtungsort seinen \emph{wahren} (vom Erdmittelpunkt aus gesehenen) Halbmesser mit $r$ und seine wahre Deklination mit $d$, für die Kulminationszeit am östlichen Ort dieselben Größen mit $\rho$ und $\delta$, dann ist

\[
t-\tau \mp \tfrac{1}{15}\left(\frac{r}{\cos d}-\frac{\rho}{\cos\delta}\right) = \Delta
\]
\noindent der wahre Rektaszensionsunterschied des Mondmittelpunkts für die beiden Beobachtungszeitpunkte. Das untere Zeichen in diesem Ausdruck gilt, wenn der westliche oder erste Mondrand beobachtet wird, das obere beim östlichen oder zweiten Mondrand. In allen Fällen, in denen die Meridiandifferenz der beiden Beobachtungsorte nicht sehr groß ist, kann man diese Korrektur ganz vernachlässigen. Auch
\SourceLeaf{449}%
sieht man leicht, dass man die absoluten Werte von $r$ und $d$ nicht sehr genau zu kennen braucht, da es hier nur auf die \emph{Änderung} dieser Größen zwischen den Beobachtungen ankommt.

Um aus $\Delta$ den Meridianunterschied der beiden Beobachtungsorte zu erhalten, benötigen wir nur ein einziges Element, nämlich die stündliche Bewegung des Mondes in AR. Bezeichnet man sie mit $h$ und die Meridiandifferenz in Zeit mit $x$, so gilt

\[
x = \left(\frac{1}{h}\,m\cdot 15\cdot 3600 - 1\right)\Delta = n\Delta
\]
\noindent wenn man den Faktor $\frac{1}{h}m\cdot15\cdot3600-1=n$ setzt. In diesem Ausdruck hängt der Faktor $m$ von der Zeitart ab, für die $h$ berechnet ist. Bedeutet $h$ die Bewegung in einer Sternzeitstunde, so ist $m=1$. Ist $h$ die Bewegung in einer mittleren Sonnenstunde, so bezeichnet $m$ das Verhältnis des mittleren Sonnentags zum Sterntag ($=1.00274$). Ist schließlich $h$ die Bewegung in einer wahren Sonnenstunde, so bezeichnet $m$ das jeweilige Verhältnis des wahren Sonnentags zum Sterntag. --- Zur Berechnung der stündlichen Bewegung $h$ ist zu bemerken, dass es nicht vorteilhaft wäre, sie durch unmittelbare Berechnung zweier nur eine Stunde auseinanderliegender Mondörter aus den Mondtafeln abzuleiten. Denn obwohl die Mondtafeln auf Zehntelsekunden berechnet sind, kann man wegen der großen Zahl der zu summierenden Gleichungen für einen daraus berechneten Mondort nie bis auf eine Sekunde bürgen. Es ist daher klar, dass die so berechnete stündliche Bewegung oft merklich fehlerhaft ausfallen würde. Weit sicherer und tatsächlich sehr genau geht man vor, wenn man mehrere Mondörter in \emph{beträchtlichen} Abständen, zum Beispiel alle 12 Stunden, berechnet und daraus durch \emph{genaue Interpolation} die stündliche Bewegung ableitet. Dadurch verteilen sich die kleinen Fehler gleichmäßig, die bei den einzelnen Örtern noch bestehen. Für dieses Verfahren kann man ohne Bedenken die in den Ephemeriden alle 12 Stunden berechneten Mondörter verwenden. --- Wegen der \emph{Veränderlichkeit} der stündlichen Bewegung selbst muss man außerdem darauf achten, das zur Bestimmung der Meridiandifferenz verwendete $h$ \emph{für das Mittel der beiden auf denselben Meridian reduzierten Beobachtungszeiten} zu berechnen. Es sei $A$ die Rektaszension des Mondes für die Stunde $\theta$ auf einem bestimmten Meridian und $A'$ dieselbe für die Stunde $\theta+1$ auf demselben Meridian. Dann ist $A'-A$ oder $h$ die stündliche Bewegung des Mondes in Rektaszension für die Zeit $\theta+30^{\rm m}$. Man berechnet also für den Meridian der Tafeln oder Ephemeriden mehrere aufeinanderfolgende stündliche Bewegungen, reduziert die beiden Beobachtungszeiten mithilfe des ungefähr bekannten Meridianunterschieds auf diesen Meridian der Tafeln oder Ephemeriden und erhält das zu verwendende $h$ für das Mittel dieser Zeiten durch Interpolation aus den stündlichen Bewegungen.

Da sich die Vergleichssterne und der Mond nicht an jedem Abend und von jedem Beobachter stets an allen Fäden des Meridianfernrohrs beobachten lassen, kann man offensichtlich nicht allen Ergebnissen dieselbe Genauigkeit beimessen. Das ist ebenfalls zu berücksichtigen. Herr Hofrat \textsc{Gauss} hatte mir die Ausdrücke für die relative Genauigkeit der einzelnen Ergebnisse schon früher freundlicherweise mitgeteilt. Sie ergeben sich leicht aus der Methode der kleinsten Quadrate. Der Mond sei an einem Ort an $l$ Fäden und am anderen an $l'$ beobachtet worden, die verschiedenen Sterne ferner an einem Ort an $a, b, c$ usw. Fäden und am anderen an $a', b', c'$ usw. Man setzt

\[
\frac{l\,l'}{l+l'}=\lambda,\qquad \frac{a\,a'}{a+a'}=\alpha,\qquad \frac{b\,b'}{b+b'}=\ObservedBeta,\qquad \frac{c\,c'}{c+c'}=\gamma,
\]
\noindent usw. Nun muss man das Mittel der Ergebnisse aus den einzelnen Vergleichssternen so bilden, dass man ihnen jeweils das Gewicht (Quadrat der Genauigkeit) $\alpha, \ObservedBeta, \gamma$ usw. zuweist. Das so erhaltene Ergebnis
\SourceLeaf{450}%
für $t-\tau$ oder für $\Delta$ hat dann das Gewicht $\dfrac{\sigma\lambda}{\sigma+\lambda}$, wenn man $\alpha+\ObservedBeta+\gamma+\text{etc.}=\sigma$ setzt. Bezeichnet man den Koeffizienten, mit dem $\Delta$ zur Bestimmung der Meridiandifferenz multipliziert wird, wie oben mit $n$, dann ist das Gewicht des Ergebnisses für diese Differenz $=\dfrac{\sigma\lambda}{(\sigma+\lambda)nn}$ und das Gewicht des Ergebnisses aus allen Bestimmungen $=\Sigma\cdot\dfrac{\sigma\lambda}{(\sigma+\lambda)nn}$. Setzt man den wahrscheinlichen Fehler eines Fadenantritts $=\varepsilon$, dann ist der wahrscheinliche Fehler des Endergebnisses

\[
=\frac{\varepsilon}{\BareRoot{\Sigma\,\dfrac{\sigma\lambda}{(\sigma+\lambda)nn}}}
\]

Wie man sieht, setzt diese Berechnung des wahrscheinlichen Fehlers voraus, dass man allen beobachteten Fadenantritten dieselbe Zuverlässigkeit beimisst, ohne zwischen Mond und Sternen, den verschiedenen Deklinationen oder den Beobachtern zu unterscheiden. Es wäre nicht schwierig, all diese Umstände zu berücksichtigen, wenn man die betreffenden Koeffizienten kennte. Doch dürfte das alles unerheblich sein. Was die verschiedenen Deklinationen betrifft, kann man ohne Bedenken alle im Tierkreis des Mondes beobachteten Fadenantritte als gleich genau betrachten, zumal die jeweiligen Begleitumstände der Beobachtung, der Zustand der Luft, der tiefere oder höhere Stand des Mondes und der Sterne usw., die Genauigkeit der Fadenantritte hier viel stärker beeinflussen werden als die größeren oder kleineren Deklinationen. Im Übrigen empfiehlt es sich, in der obigen Endformel für $s$ einen \emph{mittleren} Wert aus allen im Tierkreis des Mondes beobachteten Fadenantritten zugrunde zu legen.

\begin{center}
{\small ---\,---\,---\,---\,---\,---\,---\,---\,---\,---\,---\,---\,---\,---\,---\,---\,---\,---\,---\,---\,---\,---}
\end{center}

Diese ganze Rechnung ist, wie Sie sehen, äußerst einfach und, seit uns die Conn.\ de Tems die Mühe erspart, die Rektaszensionen des Mondes selbst zu berechnen, auch sehr bequem und kurz. --- Abschließend bemerke ich noch, dass man den Wert von $x+\Delta$ \emph{in voller Strenge} aus $\Delta$ eigentlich durch eine Reihe erhält, die nach den ungeraden Potenzen von $\Delta$ fortschreitet. Ihr erstes Glied ist $\dfrac{1}{h}m\cdot15\cdot3600\Delta$ (wobei $h$ für das Mittel der beiden Beobachtungszeiten genommen wird). Das zweite, von $\Delta^3$ abhängige Glied ist allerdings bereits so klein, dass es erst bei einer Meridiandifferenz von zwei bis zweieinhalb Stunden etwa $0^{\mathrm{s}}1$ beträgt.
Bei einem \emph{sehr großen} Längenunterschied würde ich aber ein \emph{indirektes} Verfahren vorziehen, nämlich aus einer \emph{angenommenen} Meridiandifferenz den ihr entsprechenden Unterschied der Rektaszensionen des Mondes berechnen, diesen berechneten Unterschied mit dem tatsächlich beobachteten vergleichen und aus ihrer Abweichung mithilfe der stündlichen Bewegung die Korrektur bestimmen, die an der angenommenen Meridiandifferenz noch anzubringen ist.

\begin{flushright}
\textsc{Nicolai.}
\end{flushright}
